\documentclass[11pt]{article}

\usepackage[margin=1in]{geometry}
\usepackage{amsmath,amssymb}
\usepackage{booktabs}
\usepackage[T1]{fontenc}

\title{Parallel Mean-Comparison Simulation Settings}
\author{}
\date{}

\begin{document}

\maketitle

\section{Objective}

The simulation evaluates inference for the projected null
\begin{equation}
H_0^{\rm proj}(v_1):
\left(\mu_X-\mu_Z\right)^\top v_1=0,
\label{eq:projected-null}
\end{equation}
where $v_1$ is the leading eigenvector of the common population covariance
matrix $\Sigma=\Sigma_X=\Sigma_Z$. We compare the debiased PC, oracle PC, and
plug-in PC procedures. No global-null experiment and no anchored projection
procedure are included in this simulation.

We use
\[
p\in\{100,1000\},
\qquad
N_X=N_Z\in\{50,200,400,600,800\}.
\]
For every combination of $p$ and the per-group sample size, we generate 1000
independent Monte Carlo repetitions. All tests are two-sided and use
significance level $0.05$.


\section{Population covariance matrices}

For each $p$, we first generate a sparse symmetric matrix $A$ with zero
diagonal. We sample
\[
\operatorname{round}(8p/2)=4p
\]
distinct unordered pairs of coordinates uniformly at random. For every
selected pair, its two symmetric entries in $A$ are assigned the same
independent $\operatorname{Uniform}(0.5,1.5)$ weight. Thus, each row has eight
nonzero off-diagonal entries on average, but the locations of these entries do
not follow a block or Toeplitz pattern.

The population covariance matrix is
\begin{equation}
\Sigma=I_p+0.9\frac{A}{\lVert A\rVert_{\rm op}}.
\label{eq:population-covariance}
\end{equation}
Since every eigenvalue of $A/\lVert A\rVert_{\rm op}$ is in $[-1,1]$,
\[
\lambda_{\min}(\Sigma)\geq 0.1.
\]
We retain the first generated matrix satisfying
\begin{equation}
\lambda_1(\Sigma)-\lambda_2(\Sigma)\geq 0.2.
\label{eq:population-gap-requirement}
\end{equation}
The search uses random seed $20260825$ and allows at most 500 attempts. The
accepted matrices used throughout the simulation have the following
diagnostics.

\begin{center}
\begin{tabular}{rrrrrr}
\toprule
$p$ & Nonzero pairs & $\lambda_1(\Sigma)$
& $\lambda_2(\Sigma)$ & $\lambda_1(\Sigma)-\lambda_2(\Sigma)$
& $\lambda_{\min}(\Sigma)$ \\
\midrule
100  & 400  & 1.900 & 1.553 & 0.347 & 0.445 \\
1000 & 4000 & 1.900 & 1.600 & 0.300 & 0.389 \\
\bottomrule
\end{tabular}
\end{center}

The two accepted covariance matrices are fixed across all sample sizes and
Monte Carlo repetitions. Therefore, changes across sample sizes reflect
estimation accuracy rather than changes in the underlying covariance matrix.


\section{Population mean differences}

For either mean configuration described below, observations are generated as
\[
X_i\overset{\rm IID}{\sim}
\mathcal N_p\left(\frac{\mu_X-\mu_Z}{2},\Sigma\right),
\qquad
Z_i\overset{\rm IID}{\sim}
\mathcal N_p\left(-\frac{\mu_X-\mu_Z}{2},\Sigma\right).
\]
This symmetric parameterization makes the population mean difference exactly
$\mu_X-\mu_Z$.

Let $a_1,\ldots,a_{10}$ be the indices of the ten largest entries of $v_1$
in absolute value, ordered from largest to smallest, and define the
normalization constant
\[
D=\left(\sum_{\ell=1}^{10}v_{1,a_\ell}^2\right)^{1/2}.
\]
Under the projected null, the mean difference has exactly these ten active
coordinates. For $k=1,\ldots,5$, set
\begin{align}
\left(\mu_X-\mu_Z\right)_{a_{2k-1}}
&=\frac{v_{1,a_{2k}}}{D},
&
\left(\mu_X-\mu_Z\right)_{a_{2k}}
&=-\frac{v_{1,a_{2k-1}}}{D},
\label{eq:projected-null-mean}
\end{align}
and set all remaining coordinates to zero. In words, we divide the ten
selected coordinates into five consecutive pairs. Within every pair, we swap
the two leading-PC loadings and change the sign of the second one.

Each pair makes a zero contribution to the inner product with $v_1$:
\begin{align*}
\left(\mu_X-\mu_Z\right)^\top v_1
&=
\frac{1}{D}\sum_{k=1}^5
\left(
v_{1,a_{2k-1}}v_{1,a_{2k}}
-v_{1,a_{2k}}v_{1,a_{2k-1}}
\right)
=0.
\end{align*}
Moreover,
\begin{align*}
\lVert\mu_X-\mu_Z\rVert_2^2
&=
\frac{1}{D^2}\sum_{k=1}^5
\left(v_{1,a_{2k}}^2+v_{1,a_{2k-1}}^2\right)
=1.
\end{align*}
Thus, the projected-null signal has unit Euclidean norm, is supported on ten
genes, and is exactly orthogonal to $v_1$.

For $p=100$, the actual nonzero entries, rounded to three decimals, are
\begin{align*}
\big(&\left(\mu_X-\mu_Z\right)_{78},
\left(\mu_X-\mu_Z\right)_{51},
\left(\mu_X-\mu_Z\right)_{62},
\left(\mu_X-\mu_Z\right)_{86},
\left(\mu_X-\mu_Z\right)_{33},\\
&\left(\mu_X-\mu_Z\right)_{7},
\left(\mu_X-\mu_Z\right)_{87},
\left(\mu_X-\mu_Z\right)_{84},
\left(\mu_X-\mu_Z\right)_{99},
\left(\mu_X-\mu_Z\right)_{10}\big)\\
&=(0.316,-0.442,0.315,-0.315,0.298,-0.305,
0.289,-0.294,0.272,-0.285).
\end{align*}
For $p=1000$, the corresponding entries are
\begin{align*}
\big(&\left(\mu_X-\mu_Z\right)_{992},
\left(\mu_X-\mu_Z\right)_{338},
\left(\mu_X-\mu_Z\right)_{446},
\left(\mu_X-\mu_Z\right)_{527},
\left(\mu_X-\mu_Z\right)_{35},\\
&\left(\mu_X-\mu_Z\right)_{721},
\left(\mu_X-\mu_Z\right)_{10},
\left(\mu_X-\mu_Z\right)_{20},
\left(\mu_X-\mu_Z\right)_{800},
\left(\mu_X-\mu_Z\right)_{163}\big)\\
&=(0.368,-0.379,0.320,-0.355,0.292,-0.298,
0.281,-0.291,0.277,-0.279).
\end{align*}
Thus, the two population means are different even though
\eqref{eq:projected-null} holds. Rejection probabilities under
\eqref{eq:projected-null-mean} estimate the Type I error of the projected
tests.

For the power experiment, let $S=\{a_1,\ldots,a_{10}\}$, and let $v_{1,S}$
agree with $v_1$ on $S$ and equal zero outside $S$. Starting from the mean difference in
\eqref{eq:projected-null-mean}, we add
\begin{equation}
0.25\frac{v_{1,S}}{v_1^\top v_{1,S}}.
\label{eq:projected-alternative-addition}
\end{equation}
Consequently, the mean difference under the alternative satisfies
\begin{equation}
\left(\mu_X-\mu_Z\right)^\top v_1=0.25.
\label{eq:projected-alternative}
\end{equation}
The addition in \eqref{eq:projected-alternative-addition} is supported on the
same ten coordinates as the projected-null mean difference. The complete mean
difference under the alternative is therefore also supported on ten genes.


\section{Training-fold nuisance estimators}

The plug-in and debiased procedures use five-fold cross-fitting. Each group is
randomly partitioned into five equal folds. For each fold $m\in[5]$, nuisance
quantities are estimated using only
$\mathcal D^{(-m)}$, and the resulting quantities are evaluated using the held
out observations in $\mathcal D^{(m)}$. The signs of both the population and
estimated leading eigenvectors are fixed by requiring their largest entry in
absolute value to be positive.

\subsection{Thresholded covariance estimator}

Within $\mathcal D^{(-m)}$, the observations from the two groups are centered
separately using their training-fold sample means and then pooled. Their
cross-product matrix divided by $|\mathcal D^{(-m)}|-2$ is denoted by
$\widetilde\Sigma^{(-m)}$. The subtraction of two accounts for the two
estimated group means.

The diagonal entries are retained. For $j\neq j'$, the covariance estimator is
\begin{align}
\Sigma_{jj'}^{(-m)}
={}&\widetilde\Sigma_{jj'}^{(-m)}
\mathbf 1\Bigg\{
\left|\widetilde\Sigma_{jj'}^{(-m)}\right|
\geq 1.25
\left[
\left\{
\widetilde\Sigma_{jj}^{(-m)}
\widetilde\Sigma_{j'j'}^{(-m)}
+\left(\widetilde\Sigma_{jj'}^{(-m)}\right)^2
\right\}
\frac{\log p}{|\mathcal D^{(-m)}|-2}
\right]^{1/2}
\Bigg\}.
\label{eq:thresholded-covariance}
\end{align}
If the smallest eigenvalue of the matrix in
\eqref{eq:thresholded-covariance} is less than $10^{-6}$, the minimum multiple
of $I_p$ needed to make its smallest eigenvalue equal to $10^{-6}$ is added.
This diagonal shift changes neither the eigenvectors nor the eigenvalue gaps.
The leading eigenvalue and eigenvector of the resulting matrix are denoted by
$\lambda_1^{(-m)}$ and $v_1^{(-m)}$.

This estimator is an entrywise-thresholded covariance estimator. In
particular, the revision simulation does not use \texttt{PMA::SPC} or the
spiked covariance estimator used elsewhere in the manuscript code.

\subsection{Mean-difference estimator used in the correction}

The training-fold sample means are denoted by $\mu_X^{(-m)}$ and
$\mu_Z^{(-m)}$. Before constructing the one-step correction, their difference
is hard-thresholded coordinatewise. For coordinate $j$, it is retained only
when
\begin{equation}
\left|\mu_{X,j}^{(-m)}-\mu_{Z,j}^{(-m)}\right|
\geq
1.25\sqrt{2\log p}
\left(
\frac{\widehat{\operatorname{Var}}^{(-m)}(X_j)}
{4N_X/5}
+
\frac{\widehat{\operatorname{Var}}^{(-m)}(Z_j)}
{4N_Z/5}
\right)^{1/2}.
\label{eq:thresholded-mean-difference}
\end{equation}
Here $\widehat{\operatorname{Var}}^{(-m)}(X_j)$ and
$\widehat{\operatorname{Var}}^{(-m)}(Z_j)$ are the ordinary sample variances
calculated from the two training-fold samples.
The hard-thresholded difference in
\eqref{eq:thresholded-mean-difference}, rather than the unthresholded
difference, is used in
\begin{equation}
s^{(-m)}
=
\left(\lambda_1^{(-m)}I_p-\Sigma^{(-m)}\right)^+
\left(\mu_X^{(-m)}-\mu_Z^{(-m)}\right).
\label{eq:estimated-s}
\end{equation}
The unthresholded training-fold sample means are still used to center the
observations appearing in the influence functions.


\section{Compared methods}

\subsection{Oracle PC}

The oracle procedure treats the population $v_1$ as known. It estimates
\[
\theta=\left(\mu_X-\mu_Z\right)^\top v_1
\]
by
\[
\widehat\theta_{\rm or}
=
\left(\overline X-\overline Z\right)^\top v_1.
\]
Its squared standard error is the sum of the sample variance of $X_i^\top v_1$
divided by $N_X$ and the sample variance of $Z_i^\top v_1$ divided by $N_Z$.
More explicitly,
\begin{align}
\widehat\sigma_{\rm or}^2
={}&
\frac{1}{N_X(N_X-1)}
\sum_{i=1}^{N_X}
\left\{(X_i-\overline X)^\top v_1\right\}^2
\nonumber\\
&+
\frac{1}{N_Z(N_Z-1)}
\sum_{i=1}^{N_Z}
\left\{(Z_i-\overline Z)^\top v_1\right\}^2.
\label{eq:oracle-standard-error}
\end{align}
Thus, $\widehat\sigma_{\rm or}$ is the estimated standard deviation of
$\widehat\theta_{\rm or}$, not an additional model parameter. Since $v_1$ is
an eigenvector of $\Sigma$, the population variance being estimated is
\[
\operatorname{Var}(\widehat\theta_{\rm or})
=
\lambda_1(\Sigma)\left(\frac{1}{N_X}+\frac{1}{N_Z}\right).
\]
The oracle statistic is
$T_{\rm or}=\widehat\theta_{\rm or}/\widehat\sigma_{\rm or}$.

\subsection{Plug-in PC}

For each held-out fold, the plug-in estimate is
\begin{equation}
\widehat\theta_{\rm pi}^{(m)}
=
\left(\mu_X^{(m)}-\mu_Z^{(m)}\right)^\top v_1^{(-m)}.
\label{eq:plugin-fold-estimate}
\end{equation}
The five fold estimates are averaged. If
$\widehat\sigma_{\rm pi}^{2(m)}$ is the sum of the two held-out sample
variances of the projected observations divided by their respective fold
sample sizes, the implementation uses
\begin{equation}
\widehat\theta_{\rm pi}
=\frac{1}{5}\sum_{m=1}^5\widehat\theta_{\rm pi}^{(m)},
\qquad
\widehat\sigma_{\rm pi}^2
=\frac{1}{5^2}\sum_{m=1}^5
\widehat\sigma_{\rm pi}^{2(m)},
\qquad
T_{\rm pi}=\frac{\widehat\theta_{\rm pi}}
{\widehat\sigma_{\rm pi}}.
\label{eq:plugin-statistic}
\end{equation}

\subsection{Debiased PC}

Using \eqref{eq:estimated-s}, the estimated influence functions in fold $m$
are
\begin{align}
\phi_X^{(-m)}(X)
={}&s^{(-m)\top}
\left[
\left(X-\mu_X^{(-m)}\right)
\left(X-\mu_X^{(-m)}\right)^\top
-\Sigma^{(-m)}
\right]v_1^{(-m)},
\label{eq:x-influence-function}\\
\phi_Z^{(-m)}(Z)
={}&s^{(-m)\top}
\left[
\left(Z-\mu_Z^{(-m)}\right)
\left(Z-\mu_Z^{(-m)}\right)^\top
-\Sigma^{(-m)}
\right]v_1^{(-m)}.
\label{eq:z-influence-function}
\end{align}
Because the held-out group sizes are equal, the fold-specific one-step
estimate implemented in the simulation is
\begin{align}
\widehat\theta_{\rm 1s}^{(m)}
={}&\widehat\theta_{\rm pi}^{(m)}
+\frac{1}{|\mathcal D^{(m)}|}
\left\{
\sum_{X_i\in\mathcal D^{(m)}}\phi_X^{(-m)}(X_i)
+\sum_{Z_i\in\mathcal D^{(m)}}\phi_Z^{(-m)}(Z_i)
\right\}.
\label{eq:debiased-fold-estimate}
\end{align}
For variance estimation, the corresponding held-out scores are
\begin{align*}
V_X^{(-m)}(X)
&=X^\top v_1^{(-m)}+\frac{1}{2}\phi_X^{(-m)}(X),\\
V_Z^{(-m)}(Z)
&=Z^\top v_1^{(-m)}-\frac{1}{2}\phi_Z^{(-m)}(Z).
\end{align*}
The fold variance is the sum of the two sample variances of these scores
divided by their respective held-out sample sizes. The fold estimates and
variances are then combined exactly as in \eqref{eq:plugin-statistic}, giving
$\widehat\theta_{\rm 1s}$, $\widehat\sigma_{\rm 1s}$, and
\[
T_{\rm 1s}=\frac{\widehat\theta_{\rm 1s}}
{\widehat\sigma_{\rm 1s}}.
\]

For all three methods, the two-sided $p$-value is
\[
2\Phi\left(-|T|\right),
\]
where $\Phi$ is the standard normal distribution function. Rejection occurs
when this value is at most $0.05$.


\section{Monte Carlo summaries and parallel execution}

For each combination of $p$, per-group sample size, method, and simulation
target, the reported rejection rate is the average of the 1000 rejection
indicators. If the rejection rate is $\widehat q$, its reported Monte Carlo
standard error is
\[
\left\{\frac{\widehat q(1-\widehat q)}{1000}\right\}^{1/2}.
\]
Under \eqref{eq:projected-null-mean}, this rejection rate estimates Type I
error. Under \eqref{eq:projected-alternative}, it estimates power.

For computation, ten consecutive repetitions for one method and one
$(p,N_X)$ combination form a single parallel task. The task manifest has four
columns: a unique ID, the method, the setting, and the batch number. A separate
setting table maps each setting to $p$ and $N_X$. Hence, the complete experiment
contains
\[
3\times2\times5\times100=3000
\]
parallel tasks. A task has a self-describing identifier such as
\texttt{debiased\_pc\_p1000\_n0200\_b037}. Each task writes to its own file,
so completed tasks can be validated and skipped when the pipeline is resumed.
The random seed for repetition $r$ is
\[
20260829+p\times10^6+N_X\times10^3+r.
\]
The fold assignments use this seed plus one under the projected null and this
seed plus two under the projected alternative. The parallel execution order
therefore does not affect the simulated data or the cross-fitting splits.

The settings are declared in
\texttt{config/011\_mean\_comparison\_simulation\_config.R}. Population
quantities and the task manifest are prepared by
\texttt{code/011\_prepare\_mean\_comparison\_tasks.R}; a single task is run by
\texttt{code/012\_run\_mean\_comparison\_task.R}; and the computational
interface is contained in
\texttt{src/012\_mean\_comparison\_simulation\_task.R}. The parallel
dispatcher, result collector, complete runner, and plotting script are
numbered 013 through 016, respectively.

\end{document}
